\section{Background and research direction}

Post-training changes a pretrained language model from a broad next-token predictor into
a model that follows instructions, solves target tasks, and acts through tools. The central
design choice is the learning signal. Supervised fine-tuning (SFT) supplies demonstrated
tokens; reinforcement learning (RL) judges the model's own sampled behavior; process
feedback evaluates intermediate progress; and distillation transfers a teacher's denser
token distribution. These signals differ in what behavior they can identify, the model
states on which they are available, and the cost and error of obtaining them.

The distinction becomes important for long-horizon agents. A terminal reward can identify
a successful trajectory without explaining which early action caused success, while a
demonstration gives exact actions but may not teach recovery from the model's own errors.
Process critics and on-policy distillation attempt to provide denser feedback, but add an
evaluator or teacher whose reliability must also be examined. Agent training further
requires exact provenance for model actions when tool responses, context replacement,
parallel branches, or asynchronous workers separate the optimized policy from the policy
that generated the data \citep{lightman2023verify,xu2026polar,glm5team2026}.

Model training and harness engineering are complementary. Search and verification can
extract much more from a fixed model: Jeremy Berman's ARC harness raised Claude Sonnet 3.5
from $14\%$ to $53.6\%$, but used as many as 500 candidate transforms and 31 prompts per
task \citep{kamradt2024arcprogress}. Post-training instead changes the proposal policy and
can raise its effective ceiling across harnesses. o3-preview reached $75.7\%$ on ARC-AGI-1
within the leaderboard compute limit; OpenAI links the o3 family to scaled reinforcement
learning \citep{chollet2024o3arc,openai2025o3}. This is not clean out-of-distribution
evidence because o3-preview trained on $75\%$ of ARC's public training set, although its
evaluation tasks remained private. The cautious lesson is that reasoning post-training
made ARC-like transformations more reachable, while the harness remained an important
multiplier.

This paper first establishes the minimum Transformer and SFT background needed for the
later comparison. It then develops policy-gradient RL and GRPO, followed by on-policy
distillation and a synthesis of when demonstrations, outcome rewards, process feedback,
critics, or teacher distributions are appropriate. The final question is whether an
automated loop of task synthesis, verification, failure diagnosis, data resynthesis, and
curriculum construction can turn post-training into reliable continual learning rather
than repeated capability replacement \citep{shenfeld2026selfdistillation}.

The working position is that no single objective supplies all of the information required
by long-horizon agents. Method choice should instead be justified by supervision source,
credit granularity, sampling distribution, evaluator reliability, and the need to preserve
earlier capabilities.
